\section{The Appointed Texts}
\label{sec:appointed-texts}

The Missal's texts for this Sunday are those of \latin{Dominica XXVI ``per
annum''} in the \work{Missale Romanum}, editio typica tertia (2002), which the
\work{Roman Missal}, Third Edition, renders for the dioceses of the United
States; the readings, psalm and Alleluia verse are those of no.~136 of the
\work{Lectionary for Mass}, which follows the \work{Ordo lectionum Missae} of
1981. The Latin of the Missal, the approved English of the Missal and the
approved English of the Lectionary are all under copyright, and none of them
is reproduced here. Each biblical element is given at its canonical verses in
the Douay--Rheims version as revised by Bishop Challoner, which translates the
Clementine Vulgate and numbers the psalms and the prayer of Azarias as the
Vulgate does. That English is a study text and is not the wording proclaimed at
Mass. Words printed in brackets and italics belong to the verse but not to the
appointed text. Each oration is identified by the opening words of its Latin
and described, and no English is supplied for it; the two antiphons that adapt
Scripture are described, and the verses they draw on are printed beside them as
their sources and not as their words.

\Needspace{8\baselineskip}
\subsection{Entrance Antiphon}
\label{proper-entrance-antiphon}

\properlocus{\latin{Omnia, quae fecisti nobis, Domine}. The Missal cites
Daniel 3:31, 29, 30, 43, 42 and prints no \latin{Cf.}}

\noindent A confession addressed to the Lord in the first person plural:
everything he has done to his people he did in true judgment, because they
sinned against him and did not obey his commands; he is asked to give glory to
his name and to deal with them according to the abundance of his mercy. The
antiphon adapts five verses of the prayer of Azarias, which the Vulgate and the
Douay--Rheims number within Daniel 3, and it quotes none of them. It condenses
the two parallel clauses of v.~31 into one and adds \latin{Domine}; it gives the
confession of vv.~29--30 in other words (\latin{peccavimus tibi},
\latin{mandatis tuis non oboedivimus}, where the Clementine Vulgate reads
\latin{peccavimus \ldots\ recedentes a te} and \latin{praecepta tua non
audivimus}); from v.~43 it takes only the plea for glory to God's name, and from
v.~42 only the plea for mercy, without its first clause and without
\latin{iuxta mansuetudinem tuam}. No English is printed for the antiphon. The
verses it draws on follow in canonical order.

\begin{englishwitness}{Douay--Rheims (Challoner), Dan 3:29--31, 42--43, the antiphon's source}
\textsuperscript{29}For we have sinned, and committed iniquity, departing from
thee: and we have trespassed in all things: \textsuperscript{30}And we have not
hearkened to thy commandments, nor have we observed nor done as thou hadst
commanded us, that it might go well with us. \textsuperscript{31}Wherefore, all
that thou hast brought upon us, and every thing that thou hast done to us, thou
hast done in true judgment:

\textsuperscript{42}Put us not to confusion, but deal with us according to thy
meekness, and according to the multitude of thy mercies. \textsuperscript{43}And
deliver us, according to thy wonderful works, and give glory to thy name, O
Lord:
\end{englishwitness}

\Needspace{8\baselineskip}
\subsection{Collect}
\label{proper-collect}

\properlocus{\latin{Deus, qui omnipotentiam tuam}. Long conclusion.}

\noindent The address places the chief showing of God's omnipotence in his
sparing and his mercy. The petition asks that his grace be multiplied upon the
assembly, so that those who run toward what he has promised may come to share
the goods of heaven. No English is printed.

\Needspace{8\baselineskip}
\subsection{First Reading}
\label{proper-first-reading}

\properlocus{Ezekiel 18:25--28. Lectionary no.~136.}

\begin{englishwitness}{Douay--Rheims (Challoner), Ezek 18:25--28}
\textsuperscript{25}And you have said: The way of the Lord is not right. Hear
ye, therefore, O house of Israel: Is it my way that is not right, and are not
rather your ways perverse? \textsuperscript{26}For when the just turneth himself
away from his justice, and committeth iniquity, he shall die therein: in the
injustice that he hath wrought he shall die. \textsuperscript{27}And when the
wicked turneth himself away from his wickedness, which he hath wrought, and
doeth judgment, and justice: he shall save his soul alive.
\textsuperscript{28}Because he considereth and turneth away himself from all
his iniquities which he hath wrought, he shall surely live, and not die.
\end{englishwitness}

\Needspace{8\baselineskip}
\subsection{Responsorial Psalm}
\label{proper-responsorial-psalm}

\properlocus{Psalm 25 (24):4--5, 6--7, 8--9; the response is taken from
v.~6a. Lectionary no.~136.}

\begin{englishwitness}{Douay--Rheims (Challoner), Ps 24:6, first half, as the response}
Remember, O Lord, thy bowels of compassion\notread{; and thy mercies that are
from the beginning of the world}.
\end{englishwitness}

\begin{englishwitness}{Douay--Rheims (Challoner), Ps 24:4--9}
\textsuperscript{4}\notread{Let all them be confounded that act unjust things
without cause.} Shew, O Lord, thy ways to me, and teach me thy paths.
\textsuperscript{5}Direct me in thy truth, and teach me; for thou art God my
Saviour; and on thee have I waited all the day long.\footnote{The bishops'
conference's page of readings for the day cites v.~5 whole but prints no line
answering to its last clause, on waiting for the Lord all the day long;
whether the printed Lectionary includes it is not known.}

\textsuperscript{6}Remember, O Lord, thy bowels of compassion; and thy mercies
that are from the beginning of the world. \textsuperscript{7}The sins of my
youth and my ignorances do not remember. According to thy mercy remember thou
me: for thy goodness' sake, O Lord.

\textsuperscript{8}The Lord is sweet and righteous: therefore he will give a
law to sinners in the way. \textsuperscript{9}He will guide the mild in
judgment: he will teach the meek his ways.
\end{englishwitness}

\noindent The Vulgate's v.~4 has a first clause, on the confusion of those who
act unjustly, which the Hebrew numbering places at the end of v.~3; the
Lectionary's v.~4 begins with ``Shew, O Lord, thy ways to me''.

\Needspace{8\baselineskip}
\subsection{Second Reading, Longer Form}
\label{proper-second-reading-long}

\properlocus{Philippians 2:1--11. Lectionary no.~136, the longer form.}

\begin{englishwitness}{Douay--Rheims (Challoner), Phil 2:1--11}
\textsuperscript{1}If there be therefore any consolation in Christ, if any
comfort of charity, if any society of the spirit, if any bowels of
commiseration: \textsuperscript{2}Fulfil ye my joy, that you be of one mind,
having the same charity, being of one accord, agreeing in sentiment.
\textsuperscript{3}Let nothing be done through contention: neither by vain
glory. But in humility, let each esteem others better than themselves:
\textsuperscript{4}Each one not considering the things that are his own, but
those that are other men's. \textsuperscript{5}For let this mind be in you,
which was also in Christ Jesus: \textsuperscript{6}Who being in the form of
God, thought it not robbery to be equal with God: \textsuperscript{7}But
emptied himself, taking the form of a servant, being made in the likeness of
men, and in habit found as a man. \textsuperscript{8}He humbled himself,
becoming obedient unto death, even to the death of the cross.
\textsuperscript{9}For which cause, God also hath exalted him and hath given
him a name which is above all names: \textsuperscript{10}That in the name of
Jesus every knee should bow, of those that are in heaven, on earth, and under
the earth: \textsuperscript{11}And that every tongue should confess that the
Lord Jesus Christ is in the glory of God the Father.
\end{englishwitness}

\Needspace{8\baselineskip}
\subsection{Second Reading, Shorter Form}
\label{proper-second-reading-short}

\properlocus{Philippians 2:1--5. Lectionary no.~136, the shorter form.}

\begin{englishwitness}{Douay--Rheims (Challoner), Phil 2:1--5}
\textsuperscript{1}If there be therefore any consolation in Christ, if any
comfort of charity, if any society of the spirit, if any bowels of
commiseration: \textsuperscript{2}Fulfil ye my joy, that you be of one mind,
having the same charity, being of one accord, agreeing in sentiment.
\textsuperscript{3}Let nothing be done through contention: neither by vain
glory. But in humility, let each esteem others better than themselves:
\textsuperscript{4}Each one not considering the things that are his own, but
those that are other men's. \textsuperscript{5}For let this mind be in you,
which was also in Christ Jesus:
\end{englishwitness}

\noindent The shorter form ends where the hymn of vv.~6--11 begins.

\Needspace{8\baselineskip}
\subsection{Alleluia Verse}
\label{proper-gospel-acclamation}

\properlocus{John 10:27. Lectionary no.~136. The Lectionary adds a speaker
formula, \latin{dicit Dominus} in the Latin \work{Ordo}, which the verse itself
does not have, and prints no \latin{Cf.}}

\begin{englishwitness}{Douay--Rheims (Challoner), Jn 10:27}
My sheep hear my voice. And I know them: and they follow me.
\end{englishwitness}

\noindent No English is printed for the added formula.

\Needspace{8\baselineskip}
\subsection{Gospel}
\label{proper-gospel}

\properlocus{Matthew 21:28--32. Lectionary no.~136. The Latin \work{Ordo}
opens the reading with an incipit naming the chief priests and elders of the
people, whom the narrative introduces at 21:23; the incipit is not printed
here as Scripture.}

\begin{englishwitness}{Douay--Rheims (Challoner), Mt 21:28--32}
\textsuperscript{28}But what think you? A certain man had two sons: and coming
to the first, he said: Son, go work to day in my vineyard.
\textsuperscript{29}And he answering, said: I will not. But afterwards, being
moved with repentance, he went. \textsuperscript{30}And coming to the other, he
said in like manner. And he answering said: I go, Sir. And he went not.
\textsuperscript{31}Which of the two did the father's will? They say to him:
The first. Jesus saith to them: Amen I say to you that the publicans and the
harlots shall go into the kingdom of God before you. \textsuperscript{32}For
John came to you in the way of justice: and you did not believe him. But the
publicans and the harlots believed him: but you, seeing it, did not even
afterwards repent, that you might believe him.
\end{englishwitness}

\Needspace{8\baselineskip}
\subsection{Prayer over the Offerings}
\label{proper-prayer-over-offerings}

\properlocus{\latin{Concede nobis, misericors Deus}. Short conclusion.}

\noindent The merciful God is asked to find the assembly's offering
acceptable, and through it to open to them the fountain of every blessing,
\latin{fons omnis benedictionis}. No English is printed.

\Needspace{8\baselineskip}
\subsection{Communion Antiphon, First Option}
\label{proper-communion-antiphon-a}

\properlocus{\latin{Memento verbi tui servo tuo, Domine}. The Missal cites
\latin{Cf.} Psalm 118:49--50, which is Psalm 119:49--50 in the Hebrew
numbering.}

\noindent The servant asks God to remember the word by which he gave him hope,
and says that this has been his comfort in his lowliness. The antiphon reads
\latin{Memento} where the Clementine Vulgate reads \latin{Memor esto}, adds
\latin{Domine}, and ends before the last clause of v.~50, which is printed in
brackets below. No English is printed for the antiphon.

\begin{englishwitness}{Douay--Rheims (Challoner), Ps 118:49--50, the antiphon's source}
\textsuperscript{49}Be thou mindful of thy word to thy servant, in which thou
hast given me hope. \textsuperscript{50}This hath comforted me in my
humiliation\notread{: because thy word hath enlivened me}.
\end{englishwitness}

\Needspace{8\baselineskip}
\subsection{Communion Antiphon, Second Option}
\label{proper-communion-antiphon-b}

\properlocus{\latin{In hoc cognovimus caritatem Dei}. 1~John 3:16, printed
under \latin{Vel} and without \latin{Cf.} The Missal's Latin agrees with the
Clementine Vulgate word for word.}

\begin{englishwitness}{Douay--Rheims (Challoner), 1~Jn 3:16}
In this we have known the charity of God, because he hath laid down his life
for us: and we ought to lay down our lives for the brethren.
\end{englishwitness}

\Needspace{8\baselineskip}
\subsection{Prayer after Communion}
\label{proper-prayer-after-communion}

\properlocus{\latin{Sit nobis, Domine, reparatio mentis et corporis}.
Concludes \latin{Qui vivit et regnat}.}

\noindent The heavenly mystery is asked to be the restoring of mind and body,
so that those who share the Lord's sufferings in proclaiming his death may be
his fellow heirs in glory. No English is printed.
